\section{Molecular speeds}

	\subsection{Average or mean speed}
		The average speed of the gas is defined as
		\begin{align*}
			v_\text{avg}=\expval{v}&=\dfrac{1}{N}\int\limits_{0}^{\infty}{v\dd{N_v}}
		\end{align*}
		Substituting for $\dd{N_v}$ from the Maxwellian density function, we have
		\begin{align*}
			v_\text{avg}&=\dfrac{1}{N}\times 4\pi N\qty(\dfrac{m}{2\pi k_BT})^{3/2}\underbrace{\int\limits_{0}^{\infty}{v^3e^{-mv^2/2k_BT}\dd{v}}}_{=\dfrac{1}{2\qty(m/2k_BT)^2}} \\
				&=\cancelto{2}{4}\pi\qty(\dfrac{m}{2\pi k_BT})^{3/2}\dfrac{1}{\cancel{2}\qty(m/2k_BT)^2} \\
				&=2\pi\qty(\dfrac{m}{2\pi k_BT})^{3/2}\qty(\dfrac{2k_BT}{m})^2 \\
				&=\dfrac{2\pi^{1-\sfrac{3}{2}}m^{\sfrac{3}{2}-2}}{2^{\sfrac{3}{2}-2}k_B^{\sfrac{3}{2}-2}T^{\sfrac{3}{2}-2}}=\dfrac{2^{\sfrac{3}{2}}k_B^{\sfrac{1}{2}}T^{\sfrac{1}{2}}}{\pi^{\sfrac{1}{2}}m^{\sfrac{1}{2}}} \\
				&=\qty(\dfrac{2^3k_BT}{\pi m})^{\sfrac{1}{2}} \\
			\implies\Aboxed{v_\text{avg}&=\sqrt{\dfrac{8k_BT}{\pi m}}}
		\end{align*}
	
	\subsection{Root mean squared (RMS) speed}
		The root mean squared (RMS) speed is defined as
		\begin{align*}
			v_\text{rms}&=\sqrt{\expval{v^2}}
		\end{align*}
		
		Now, we already have that
		\begin{align*}
			\expval{v^2}=\dfrac{3}{2B}\qq{and}B=\dfrac{m}{2k_BT}
		\end{align*}
		The above give us
		\begin{align*}
			\expval{v^2}&=\dfrac{3}{\cancel{2}}\times\dfrac{\cancel{2}k_BT}{m}=\dfrac{3k_BT}{m}
		\end{align*}
		
		Therefore the RMS speed is
		\begin{align*}
			\Aboxed{v_\text{rms}&=\sqrt{\dfrac{3k_BT}{m}}}
		\end{align*}
	
	\subsection{Most probable speed}
		\setcounter{equation}{0}
		The most probable speed $v_p$ is defined as the speed at which the probability of finding a molecule is maximum.
		\begin{itemize}
			\item This is obtained by finding the root of the equation $\dv{f_v}{v}=0$.
		\end{itemize}
		
		We have
		\begin{align}
			f_v&=4\pi\qty(\dfrac{m}{2\pi k_BT})^{3/2}v^2e^{-mv^2/2k_BT} \label{eqn:mps-fv}
		\end{align}
		
		Differentiating equation (\ref{eqn:mps-fv}) with respect to $v$ and equating to 0 gives $v=v_p$, we have
		\begin{align*}
			\dv{f_v}{v}&=0 \\
			\implies 4\pi\qty(\dfrac{m}{2\pi k_BT})^{3/2}\dv{v}\qty[v^2e^{-mv^2/2k_BT}]&=0 \\
			\implies\dv{v}\qty[v^2e^{-mv^2/2k_BT}]&=0 \\
			\implies 2v_pe^{-mv^2/2k_BT}+v_p^2\qty(-\dfrac{m}{2k_BT})2v_pe^{-mv^2/2k_BT}&=0 \\
			\implies 2v_pe^{-mv^2/2k_BT}\qty[1-\dfrac{mv_p^2}{2k_BT}]&=0
		\end{align*}
		On the LHS of the above, the factor $e^{-mv^2/2k_BT}$ cannot be zero for any finite value of $v$. Hence we must have
		\begin{align*}
			1-\dfrac{mv_p^2}{2k_BT}&=0 \\
			\implies\dfrac{mv_p^2}{2k_BT}&=1 \\
			\implies v_p^2&=\dfrac{2k_BT}{m} \\
			\implies\Aboxed{v_p&=\sqrt{\dfrac{2k_BT}{m}}}
		\end{align*}
		
	\subsection*{Summary}
		The three molecular speeds are as follows:
		\begin{align*}
			v_\text{avg}=\sqrt{\dfrac{8k_BT}{\pi m}},\qquad v_\text{rms}=\sqrt{\dfrac{3k_BT}{m}},\qquad v_p=\sqrt{\dfrac{2k_BT}{m}}
		\end{align*}
		
		This gives their ratios to be
		\begin{align*}
			v_\text{avg}:v_\text{rms}:v_p=\sqrt{\dfrac{8}{\pi}}:\sqrt{3}:\sqrt{2}\approx 1.6:1.7:1.4
		\end{align*}
		The above shows that
		\begin{align*}
			v_p<v_\text{avg}<v_\text{rms}
		\end{align*}
		\begin{figure}[h!]
			\centering
			\includegraphics*[width=0.8\textwidth]{figures/fig_max-mol-speeds.png}
		\end{figure}
		
		\begin{table*}[h!]
			\centering
			\caption{Table showing values of $v_\text{avg}$, $v_\text{rms}$ and $v_p$ for different gases at STP}
			\begin{tabular}{cccc}
				\toprule[1.5pt]
				\textbf{Gas} & {$\boldsymbol{v_\text{avg}\;(\unit{\metre\per\second})}$} & {$\boldsymbol{v_\text{rms}\;(\unit{\metre\per\second})}$} & {$\boldsymbol{v_p\;(\unit{\metre\per\second})}$} \\
				\midrule
				{\ce{H2}} & {1695} & {1838} & {1501} \\
				{\ce{H2O}} & {567} & {615} & {502} \\
				{\ce{N2}} & {455} & {493} & {403} \\
				{Air} & {447} & {485} & {396} \\
				{\ce{O2}} & {425} & {461} & {376} \\
				{\ce{CO2}} & {362} & {393} & {321} \\
				\bottomrule[1.5pt]
			\end{tabular}
		\end{table*}
		
	\subsection*{Exercises}
		\begin{enumerate}
			\item For a Maxwellian gas, show that $\expval{v}\expval{\dfrac{1}{v}}=\dfrac{4}{\pi}$.
			
			\item Calculate the most probable speed, average speed and root mean square speed for \ce{O2} molecules at $300\,\unit{\kelvin}$ using the following data:\vspace{-1.5ex}
			\begin{align*}
				m(\ce{O2})=\num{5.31e-26}\,\unit{\kilo\gram}\qq{and}k_B=\num{1.38e-23}\,\unit{\joule\per\kelvin}
			\end{align*}
			
			\item Calculate the fraction of molecules of a gas within $1\%$ of the most probable speed at STP. Will this fraction remain the same at all temperatures?
			
			\item Calculate the temperature at which\vspace{-1.5ex}
			\begin{enumerate}[(a)]
				\item Root mean square speed of \ce{O2} molecules exceeds their most probable speed by $150\;\unit{\metre\per\second}$.
				\item The molecular velocity distribution function for \ce{O2} peaks for $400\;\unit{\metre\per\second}$.
			\end{enumerate}
		\end{enumerate}
		
% Link for online Matplotlib plotting: https://matplotlib.codeutility.io/